\documentclass[11pt,a4paper]{article}
\usepackage[utf8]{inputenc}
\usepackage[T1]{fontenc}
\usepackage[czech]{babel}
\usepackage[margin=1.9cm,top=2.0cm,bottom=2.0cm]{geometry}
\usepackage{amsmath,amssymb}
\usepackage{parskip}
\usepackage{enumitem}
\setlist{nosep,leftmargin=1.4em}
\usepackage{booktabs}
\usepackage{array}
\usepackage{xcolor}
\definecolor{linkblue}{RGB}{20,60,150}
\definecolor{upcol}{RGB}{0,110,60}
\definecolor{downcol}{RGB}{160,50,0}
\definecolor{markcol}{RGB}{176,84,0}
\definecolor{markbg}{RGB}{253,246,236}
\usepackage[most]{tcolorbox}
\newtcolorbox{poznbox}[1][]{breakable,colback=markbg,colframe=markcol,boxrule=0.6pt,
  arc=1.5mm,left=2mm,right=2mm,top=1mm,bottom=1mm,title={#1},fonttitle=\bfseries}
\usepackage[colorlinks=true,linkcolor=linkblue,urlcolor=linkblue]{hyperref}

\setlength{\parskip}{0.35em}

% směr: co je lepší
\newcommand{\up}{\textcolor{upcol}{$\boldsymbol{\uparrow}$}}
\newcommand{\down}{\textcolor{downcol}{$\boldsymbol{\downarrow}$}}
\newcommand{\neut}{$\circ$}

% tabulka měr
\newenvironment{mirytab}{%
  \par\medskip\noindent\small
  \begin{tabular}{@{}>{\raggedright\arraybackslash}p{3.15cm}
                    >{\centering\arraybackslash}p{0.85cm}
                    >{\raggedright\arraybackslash}p{5.3cm}
                    >{\raggedright\arraybackslash}p{6.1cm}@{}}
  \toprule
  \textbf{Míra} & & \textbf{Co měří} & \textbf{Na co pozor} \\
  \midrule}%
  {\bottomrule\end{tabular}\par\medskip}

\setcounter{tocdepth}{1}

\title{\textbf{Míry a~metriky ke~SZZ}\\[2mm]
\large Co která měří a~kterým směrem je lépe}
\author{SZZ Umělá inteligence, MFF UK}
\date{srpen 2026}

\begin{document}
\maketitle

\begin{center}
\begin{minipage}{0.93\textwidth}
\small
Přehled všech měr, kvalit a~kritérií, které se v~okruzích objevují --- \textbf{bez vzorců}.
U~každé jen to, \emph{co popisuje}, \emph{kterým směrem je lépe} a~\emph{kde je past}.
Vzorce jsou ve \texttt{04-vzorce}, výklad ve \texttt{01}--\texttt{03}.
\end{minipage}
\end{center}

\begin{center}
\begin{minipage}{0.93\textwidth}
\begin{poznbox}[Legenda směru]
\small
\up~= \textbf{větší je lepší} \quad\textbullet\quad
\down~= \textbf{menší je lepší} \quad\textbullet\quad
\neut~= \textbf{deskriptivní} --- popisuje vlastnost, ne kvalitu; \uv{lepší} u~ní nedává
smysl (vzdálenost, hustota sítě, reciprocita).

\textbf{Pozor na past se směrem:} u~několika měr závisí žádoucí směr na tom, \emph{co
porovnáváš}. Korelace mezi dvěma \emph{vstupy} vysoká = redundance (špatně), mezi vstupem
a~\emph{cílem} vysoká = relevance (dobře). Táž míra, opačné čtení.
\end{poznbox}
\end{minipage}
\end{center}

\tableofcontents

% =====================================================================
\section{Vyhodnocování klasifikace}
% =====================================================================

\begin{mirytab}
Accuracy \newline (správnost) & \up &
Podíl správně klasifikovaných ze všech testovaných. &
\textbf{Zavádějící na nevyvážených datech}: při 1~\% podvodů má klasifikátor
\uv{nikdy podvod} accuracy 99~\%. \\[2pt]

Chybovost \newline (\emph{error rate}) & \down &
Doplněk accuracy do 1. &
Předpokládá, že každá chyba stojí stejně --- což skoro nikdy neplatí. \\[2pt]

Precision & \up &
Z~vzorů, které model \textbf{označil} za pozitivní, kolik jich pozitivních opravdu je. &
Hodnotí \emph{jen pozitivní třídu}. Jde snadno vyhnat k~1 tím, že model označí jen
pár jistých případů --- proto nikdy sama. \\[2pt]

Recall \newline (TPR, senzitivita) & \up &
Ze \textbf{skutečně} pozitivních vzorů kolik jich model našel. &
Taky jen pozitivní třída. Jde vyhnat k~1 tím, že model označí všechno --- proto vždy
spolu s~precision. \\[2pt]

Specificita (TNR) & \up &
Ze skutečně negativních kolik jich model správně nechal být. &
Doplněk FPR do 1. Druhá osa ROC křivky. \\[2pt]

FPR \newline (falešně pozitivní) & \down &
Podíl skutečně negativních, které model označil za pozitivní. &
\uv{Falešný poplach}. U~spamu nebo lékařské diagnózy bývá dražší než FN. \\[2pt]

$F_1$-míra & \up &
Kombinuje precision a~recall do jednoho čísla (harmonický průměr). &
Harmonický průměr je blíž tomu \textbf{menšímu} z~obou --- $F_1$ je velká, jen když jsou
velké \emph{obě}. Proto neošidíš. \\[2pt]

AUC \newline (plocha pod ROC) & \up &
Pravděpodobnost, že náhodný pozitivní vzor dostane vyšší skóre než náhodný negativní. &
$0{,}5$ = náhoda, $1{,}0$ = dokonalé. \textbf{Nezávisí na prahu ani na poměru tříd} ---
proto se hodí tam, kde accuracy klame. \\[2pt]

Lift (lift analýza) & \up &
Kolikrát je model lepší než náhodné oslovení --- kolik pozitivních je v~nejlépe
oskórované přihrádce oproti průměru. &
Marketingová míra: rozdíl mezi lift křivkou a~diagonálou je přínos modelu.
Rozhodnutí \uv{kolik zásilek} pak plyne z~ekonomiky, ne z~míry. \\[2pt]

Střední hodnota ceny & \down &
Průměrná cena chyb, když má každý typ záměny jinou cenu. &
Zobecnění chybovosti (ta je speciální případ s~cenou 0/1). Realističtější:
nedoručený ham vadí víc než nerozpoznaný spam. \\
\end{mirytab}

\begin{poznbox}[Jak si zapamatovat precision vs.\ recall]
\small
\textbf{Precision} se ptá na to, co model \emph{řekl}: \uv{kolik z~mých obvinění bylo
oprávněných?} \quad
\textbf{Recall} se ptá na to, co ve světě \emph{je}: \uv{kolik viníků jsem dopadl?}

Jmenovatel to prozradí: u~precision je v~něm \emph{predikce} modelu, u~recallu
\emph{skutečnost}. Obě přitom hodnotí \textbf{pouze pozitivní třídu} --- o~tom, jak dobře
model rozpoznává negativní, neříkají nic (na to je specificita).
\end{poznbox}

% =====================================================================
\section{Pravidla a~asociace}
% =====================================================================

\begin{mirytab}
Support (podpora) & \up &
Jak \emph{často} je pravidlo použitelné --- podíl transakcí obsahujících celou kombinaci. &
Hlídá statistickou váhu. Bez něj projde pravidlo s~dokonalou konfidencí postavené
na třech transakcích. \\[2pt]

Confidence \newline (spolehlivost) & \up &
Jak \emph{spolehlivě} pravidlo platí --- z~transakcí s~předpokladem, kolik jich má
i~závěr. &
Je to \emph{precision} pravidla. Sama nestačí: vysoká confidence může jen kopírovat to,
že závěr je obecně častý. \\[2pt]

Lift (zdvih) & \up &
Kolikrát častěji platí závěr při splněném předpokladu, než by odpovídalo náhodě. &
\textbf{Referenční hodnota je 1}, ne 0. $>1$ souvislost, $=1$ nezávislost,
$<1$ horší než náhoda --- tam dá lepší pravidlo \emph{negace} závěru. \\[2pt]

t-weight \newline (typičnost) & \up &
Jakou část cílové třídy daný zobecněný záznam pokrývá. &
Je to \emph{recall} pravidla vůči cílové třídě. Součet přes všechny záznamy je 100~\%. \\[2pt]

d-weight \newline (rozlišovací váha) & \up &
Z~objektů pokrytých záznamem kolik jich patří do cílové třídy. &
Je to \emph{precision} pravidla. Vždy srovnat s~\textbf{apriorním} podílem třídy ---
d-weight 0,6 při třídě zabírající 60~\% dat nerozlišuje nic. \\[2pt]

MIS \newline (minimum item support) & \neut &
Individuální práh supportu pro jednotlivou položku v~MS-Apriori. &
Zadává uživatel podle vzácnosti položky. Práh množiny = \emph{minimum} MIS jejích položek
--- proto přestane platit uzavřenost dolů. \\
\end{mirytab}

% =====================================================================
\section{Relevance atributů a~dělení uzlů}
% =====================================================================

\begin{mirytab}
$\chi^2$ & \up &
Jak moc se pozorované četnosti liší od těch očekávaných \emph{při nezávislosti}. &
Použitelné jen při dost velkých četnostech (očekávaná $\ge 5$); jinak Fisherův exaktní
test. Jako test má i~kritickou hodnotu, nejen pořadí. \\[2pt]

Entropie $H(A)$ & \down &
Kolik \textbf{nejistoty o~třídě zůstane} po rozdělení podle atributu (vážený průměr
přes větve). &
Vážení je podstatné --- brání tomu, aby vyhrávaly atributy s~mnoha titěrnými čistými
větvemi. \\[2pt]

Informační zisk & \up &
Kolik bitů nejistoty o~třídě dělení \textbf{ubralo}. &
Totéž rozhodnutí jako minimalizace $H(A)$ (entropie rodiče je konstanta).
\textbf{Nadržuje atributům s~mnoha hodnotami} --- v~limitě vyhraje rodné číslo. \\[2pt]

Poměrný informační zisk & \up &
Informační zisk normalizovaný tím, jak moc atribut větví. &
Oprava předchozí vady (C4.5). Jmenovatel penalizuje nadměrné větvení. \\[2pt]

Gini index & \down &
Nečistota uzlu --- pravděpodobnost, že dva náhodně tažené prvky jsou z~různých tříd. &
Stejný princip jako entropie, jiná stupnice (max $0{,}5$ pro 2 třídy vs.\ 1 u~entropie),
takže \textbf{hodnoty spolu neporovnávej}. Levnější (bez logaritmů). \\[2pt]

Kritérium $\Phi$ (CART) & \up &
Jak dobře binární dělení oddělí třídy mezi levý a~pravý podstrom. &
Maxima nabývá, když jsou všechny vzory dané třídy v~témže podstromu. \\[2pt]

Vzájemná informace & \up &
Kolik informace o~třídě atribut nese. &
Systematicky zvýhodňuje atributy s~mnoha hodnotami --- proto se normalizuje. \\[2pt]

Informační míra závislosti & \up &
Vzájemná informace normalizovaná entropií cílového atributu. &
Díky normalizaci \textbf{srovnatelná napříč úlohami}, na rozdíl od samotné MI. \\[2pt]

Pearsonovo $r$ & \up/\down &
Sílu a~směr \textbf{lineárního} vztahu dvou numerických veličin. &
\textbf{Směr závisí na tom, co porovnáváš}: vysoké $|r|$ mezi vstupem a~cílem = relevance
(dobře), mezi dvěma vstupy = redundance (jeden zahodit). $r=0$ \emph{neznamená}
nezávislost. \\[2pt]

Spearmanovo $\rho$ & \up/\down &
Totéž, ale pro \textbf{libovolný monotónní} vztah (počítá se z~pořadí). &
Robustní vůči odlehlým hodnotám, invariantní vůči monotónní transformaci, funguje
i~na ordinálních atributech. \\
\end{mirytab}

% =====================================================================
\section{Vzdálenosti a~shlukování}
% =====================================================================

\begin{mirytab}
Euklidovská, Manhattan, Čebyševova & \neut &
Vzdálenost dvou vzorů (speciální případy Minkowského metriky). &
\textbf{Není to míra kvality} --- menší znamená podobnější, ne \uv{lepší}. Předpokládají
stejný rozptyl všech veličin $\Rightarrow$ nutná normalizace. \\[2pt]

Mahalanobisova vzdálenost & \neut &
Vzdálenost \uv{po narovnání} prostoru podle rozptylů \emph{a~korelací} dat. &
Zahrne to, co euklidovská nevidí. Vyžaduje invertovatelnou kovarianční matici;
na malých datech je její odhad zašuměný. \\[2pt]

Účelová funkce \newline $k$-means / $k$-medoids & \down &
Součet vzdáleností bodů k~jejich nejbližším reprezentantům. &
Klesá s~rostoucím $k$ vždy --- \textbf{nelze jí volit $k$}. Na to slouží loket, siluety
či dendrogram. \\[2pt]

Vzdálenost klastrů (\emph{linkage}) & \neut &
Jak se měří vzdálenost mezi dvěma skupinami při hierarchickém shlukování ---
single, complete, average, centroidní. &
\textbf{Volba určuje tvar} nalezených klastrů: single řetězí a~zvládne protáhlé tvary,
complete dělá kompaktní kulovité. \\
\end{mirytab}

% =====================================================================
\section{Analýza sociálních sítí}
% =====================================================================

\begin{mirytab}
Stupňová centralita & \up &
Kolik lidí uzel dosáhne \textbf{přímo}; popularita a~přímý vliv. &
Nejlevnější, ale nejkrátkozrakejší --- nevidí, kde uzel v~síti leží. \\[2pt]

Mezilehlá centralita & \up &
Jak často uzel leží na nejkratších cestách mezi jinými. &
Označuje \uv{mosty}: místa, jejichž odstraněním se síť \textbf{rozpadne}. \\[2pt]

Blízkostní centralita & \up &
Jak rychle se od uzlu dostane informace ke všem ostatním. &
Převrácená průměrná vzdálenost. Užitečná, jde-li o~rychlost šíření. \\[2pt]

Vlastní centralita & \up &
Jak dobře je uzel spojen s~\emph{jinými dobře propojenými} uzly (rekurzivně). &
Blízká příbuzná PageRanku. Nutná normalizace, jinak nekonverguje. \\[2pt]

Překryv sousedství & \up &
Sílu vazby --- jak velký podíl přátel mají oba konce společný. &
Malý překryv = \emph{most}: slabá vazba, ale zprostředkuje nové informace
(síla slabých vazeb). \\[2pt]

Hustota sítě & \neut &
Podíl existujících hran na počtu možných. &
Deskriptivní. Klika má 1. Orientované grafy mají poloviční hustotu, protože možných
hran je dvakrát víc. \\[2pt]

Klastrovací koeficient & \up &
Jak moc jsou sousedé uzlu propojení mezi sebou (tranzitivita). &
Vysoká hodnota indikuje \textbf{komunitní strukturu}; jedna ze šesti charakteristik
reálných komplexních sítí. \\[2pt]

Reciprocita & \neut &
Podíl obousměrných vztahů. &
\textbf{Má smysl jen v~orientovaných grafech.} \\[2pt]

Modularita $Q$ & \up &
Jak moc je hran \emph{uvnitř} komunit víc, než by odpovídalo náhodě. &
Kritérium, které maximalizuje Louvain. Má \emph{rozlišovací limit} --- malé komunity
ve velkých sítích nenajde. \\[2pt]

RatioCut / NCut & \down &
Cenu rozdělení sítě, normalizovanou velikostí komunit. &
Normalizace je tam proto, že prostý minimální řez vrací nevyvážená rozdělení
(často singleton). \\[2pt]

CN, Jaccard, Adamic--Adar & \up &
Jak pravděpodobně mezi dvěma uzly vznikne vazba (sousedské míry). &
Každá opravuje vadu předchozí: CN nevidí stupně, Jaccard nevidí stupeň
\emph{prostředního} souseda, AA ho váží. \\[2pt]

Katzova míra & \up &
Nepřímou propojenost přes procházky všech délek, tlumené délkou. &
Pomáhá ve \textbf{velmi řídkých} sítích, kde společných sousedů je málo a~sousedské
míry nerozliší. \\
\end{mirytab}

% =====================================================================
\section{Evoluční a~rojové algoritmy}
% =====================================================================

\begin{mirytab}
Fitness & \up &
Kvalitu jedince vůči úloze; řídí selekci. &
U~minimalizace se otáčí. V~koevoluci není absolutní, ale \emph{kontextová} --- závisí na
ostatních jedincích. \\[2pt]

Selekční tlak & \neut &
Jak moc jsou lepší jedinci zvýhodněni. &
\textbf{Není lepší velký ani malý}: velký urychlí konvergenci, ale zabije diverzitu;
malý naopak. To je celé dilema explorace vs.\ exploatace. \\[2pt]

Doba převzetí $\tau^*$ & \neut &
Za kolik generací zaplní kopie nejlepšího jedince celou populaci. &
Kvantifikuje selekční tlak --- \emph{kratší doba} znamená \emph{vyšší tlak}.
Nezaměňovat: menší $\tau^*$ není \uv{lepší}. \\[2pt]

Řád a~definující délka schématu & \down &
Jak je schéma \uv{křehké} vůči mutaci ($o(S)$) a~křížení ($d(S)$). &
Věta o~schématech: množí se schémata \emph{krátká}, \emph{nízkého řádu}
a~\emph{nadprůměrná} --- proto u~obou menší lépe přežívá. \\[2pt]

Crowding distance & \up &
Jak řídce je okolí jedince obsazené na Pareto frontě. &
Velká hodnota = \textbf{chceme si ho nechat} (drží frontu rozprostřenou).
Krajní body dostanou $\infty$. \\[2pt]

Hypervolume & \up &
Objem prostoru kritérií, který fronta \uv{ukrojila} vůči referenčnímu bodu. &
Měří konvergenci \emph{i} rozprostření jedním číslem a~\textbf{nepotřebuje znát
skutečnou frontu}. Cena: volba referenčního bodu, exponenciální složitost. \\[2pt]

GD (generational distance) & \down &
Průměrnou vzdálenost nalezených bodů od skutečné Pareto fronty. &
Měří jen konvergenci, ne pokrytí. \textbf{Vyžaduje znát skutečnou frontu} --- na reálné
úloze tedy nepoužitelná. \\[2pt]

IGD & \down &
Vzdálenost od skutečné fronty k~nalezené (opačný směr než GD). &
Měří konvergenci \emph{i} pokrytí. Taky potřebuje skutečnou frontu. \\[2pt]

Spread & \down &
Jak nerovnoměrně jsou řešení po frontě rozložená. &
Doplněk k~hypervolume: hlídá rovnoměrnost, ne objem. \\[2pt]

Míra podobnosti $\delta$ (NEAT) & \neut &
Jak jsou si dva genomy strukturně a~parametricky vzdálené. &
Je to \emph{vzdálenost}, ne kvalita. Pod prahem $\Rightarrow$ tentýž druh, uvnitř
druhu se sdílí fitness (ochrana nových topologií). \\
\end{mirytab}

% =====================================================================
\section{Multiagentní systémy}
% =====================================================================

\begin{mirytab}
Utilita & \up &
Jak je pro agenta daný stav či běh žádoucí; cíl se zadává \emph{nepřímo} přes ni. &
Utilita běhu je lepší než utilita stavu (méně krátkozraká). Definice optimálního agenta
je elegantní, ale \textbf{nekonstruktivní}. \\[2pt]

$f = g + h$ (A*) & \down &
Odhad ceny nejlepší cesty vedoucí přes daný vrchol; expanduje se nejmenší. &
Heuristika musí být \emph{přípustná} (nenadhodnocuje), aby byl A* optimální;
konzistence je silnější a~přípustnost z~ní plyne. \\[2pt]

$Q(s,a)$ & \up &
Očekávaný diskontovaný součet budoucích odměn za akci $a$ ve stavu $s$. &
Politika se z~$Q$ získá jako argmax --- \textbf{bez znalosti modelu}, na rozdíl
od~$V$. \\[2pt]

Risk (Zeuthen) & \down &
Ochotu agenta riskovat konflikt při vyjednávání. &
\textbf{Ustupuje ten s~nižším Riskem} --- tady je menší hodnota \uv{horší pozice},
ne lepší výsledek. \\[2pt]

Sociální blahobyt & \up &
Součet utilit všech agentů --- kvalita výsledku pro skupinu. &
Nezaměňovat s~\emph{Paretovou optimalitou}: ta neříká nic o~součtu, jen že nelze nikomu
přilepšit bez uškození jinému. \\
\end{mirytab}

\begin{poznbox}[Nejčastější past: směr není vlastnost míry, ale otázky]
\small
Několik měr nemá pevný \uv{lepší} směr a~zkoušející se na to rád chytá:

\begin{itemize}
\item \textbf{Korelace} --- vysoká mezi vstupem a~cílem je \emph{relevance} (chceme),
      vysoká mezi dvěma vstupy je \emph{redundance} (nechceme).
\item \textbf{Selekční tlak a~doba převzetí} --- ani velký, ani malý není lepší; je to
      nastavitelný kompromis mezi explorací a~exploatací.
\item \textbf{Vzdálenosti} --- menší znamená podobnější, ne kvalitnější.
\item \textbf{Risk u~Zeuthena} --- nižší hodnota znamená, že agent ustoupí, tedy horší
      vyjednávací pozici.
\item \textbf{Lift} --- referenční hodnota je \textbf{1}, ne 0; pod jednou se vyplatí
      závěr pravidla znegovat.
\end{itemize}
\end{poznbox}

\end{document}
